% ============================================================
%  capitoli/operazioni_aritmetiche.tex — Matemagica
%  Le operazioni aritmetiche basilari e le loro proprietà
% ============================================================

\chapter{Operazioni aritmetiche --- Proprietà}

% --- Obiettivi di apprendimento --------------------------------
\section*{Obiettivi di apprendimento}
\begin{itemize}
  \item Conoscere le quattro operazioni aritmetiche basilari: addizione, sottrazione, moltiplicazione e divisione.
  \item Saper applicare correttamente le proprietà di ciascuna operazione.
  \item Riconoscere e prevenire gli errori più comuni nel calcolo.
  \item Comprendere la precedenza tra le operazioni.
\end{itemize}

\bigskip

% ============================================================
%  SEZIONE 1 — ADDIZIONE
% ============================================================
\section{L'addizione}

\begin{definizione}
  L'\textbf{addizione} è l'operazione che combina due o più numeri, detti \textbf{addendi}, per ottenere un unico numero chiamato \textbf{somma}.

  \[
    a + b = s
  \]

  dove $a$ e $b$ sono gli addendi e $s$ è la somma.
\end{definizione}

\subsection{Proprietà dell'addizione}

\begin{regola}
  \textbf{Proprietà commutativa}

  L'ordine degli addendi non cambia la somma:
  \[
    a + b = b + a
  \]
\end{regola}

\begin{esempio}
  \[
    3 + 7 = 7 + 3 = 10
  \]
  Cambiare l'ordine degli addendi non modifica il risultato.
\end{esempio}

\begin{regola}
  \textbf{Proprietà associativa}

  Quando si sommano tre o più addendi, il risultato non cambia raggruppandoli in modo diverso:
  \[
    (a + b) + c = a + (b + c)
  \]
\end{regola}

\begin{esempio}
  \[
    (2 + 5) + 3 = 7 + 3 = 10
  \]
  \[
    2 + (5 + 3) = 2 + 8 = 10
  \]
  Il risultato è lo stesso indipendentemente da come si raggruppano gli addendi.
\end{esempio}

\begin{regola}
  \textbf{Proprietà dell'elemento neutro}

  Aggiungere zero a qualsiasi numero lascia il numero invariato:
  \[
    a + 0 = 0 + a = a
  \]
  Lo \textbf{zero} è l'\textbf{elemento neutro} dell'addizione.
\end{regola}

\begin{esempio}
  \[
    15 + 0 = 15 \qquad \qquad 0 + 8 = 8
  \]
\end{esempio}

% ============================================================
%  SEZIONE 2 — SOTTRAZIONE
% ============================================================
\section{La sottrazione}

\begin{definizione}
  La \textbf{sottrazione} è l'operazione che, dati due numeri, calcola la loro \textbf{differenza}. Il primo numero si chiama \textbf{minuendo}, il secondo \textbf{sottraendo}.

  \[
    a - b = d
  \]

  dove $a$ è il minuendo, $b$ è il sottraendo e $d$ è la differenza.
\end{definizione}

\subsection{Proprietà della sottrazione}

\begin{regola}
  \textbf{Proprietà invariantiva}

  La differenza tra due numeri non cambia se si aggiunge (o si sottrae) uno stesso numero ad entrambi i termini:
  \[
    a - b = (a + c) - (b + c) = (a - c) - (b - c)
  \]
\end{regola}

\begin{esempio}
  \[
    18 - 6 = 12
  \]
  \[
    (18 + 4) - (6 + 4) = 22 - 10 = 12
  \]
  \[
    (18 - 3) - (6 - 3) = 15 - 3 = 12
  \]
  La differenza rimane sempre $12$.
\end{esempio}

\begin{attenzione}
  \textbf{Attenzione — La sottrazione non è commutativa!}

  $a - b \neq b - a$ in generale.

  \medskip
  Ad esempio: $10 - 3 = 7$, ma $3 - 10 \neq 7$.

  Non si può invertire l'ordine dei termini di una sottrazione senza cambiare il risultato.
\end{attenzione}

% ============================================================
%  SEZIONE 3 — MOLTIPLICAZIONE
% ============================================================
\section{La moltiplicazione}

\begin{definizione}
  La \textbf{moltiplicazione} è l'operazione che calcola il \textbf{prodotto} di due numeri, detti \textbf{fattori}.

  \[
    a \cdot b = p
  \]

  dove $a$ e $b$ sono i fattori e $p$ è il prodotto.

  \medskip
  Il simbolo utilizzato in questo libro è il \textbf{punto moltiplicatore} ($\cdot$), preferito alla croce ($\times$) per evitare confusione con la lettera $x$.
\end{definizione}

\subsection{Proprietà della moltiplicazione}

\begin{regola}
  \textbf{Proprietà commutativa}

  L'ordine dei fattori non cambia il prodotto:
  \[
    a \cdot b = b \cdot a
  \]
\end{regola}

\begin{esempio}
  \[
    4 \cdot 6 = 6 \cdot 4 = 24
  \]
\end{esempio}

\begin{regola}
  \textbf{Proprietà associativa}

  Raggruppare i fattori in modo diverso non cambia il prodotto:
  \[
    (a \cdot b) \cdot c = a \cdot (b \cdot c)
  \]
\end{regola}

\begin{esempio}
  \[
    (2 \cdot 3) \cdot 5 = 6 \cdot 5 = 30
  \]
  \[
    2 \cdot (3 \cdot 5) = 2 \cdot 15 = 30
  \]
\end{esempio}

\begin{regola}
  \textbf{Proprietà distributiva rispetto all'addizione e alla sottrazione}

  La moltiplicazione si distribuisce sull'addizione e sulla sottrazione:
  \[
    a \cdot (b + c) = a \cdot b + a \cdot c
  \]
  \[
    a \cdot (b - c) = a \cdot b - a \cdot c
  \]
\end{regola}

\begin{esempio}
  \[
    3 \cdot (4 + 5) = 3 \cdot 4 + 3 \cdot 5 = 12 + 15 = 27
  \]
  Verifica: \quad $3 \cdot (4 + 5) = 3 \cdot 9 = 27$ \quad \textit{(corretto)}

  \medskip
  \[
    5 \cdot (8 - 3) = 5 \cdot 8 - 5 \cdot 3 = 40 - 15 = 25
  \]
  Verifica: \quad $5 \cdot (8 - 3) = 5 \cdot 5 = 25$ \quad \textit{(corretto)}
\end{esempio}

\begin{regola}
  \textbf{Proprietà dell'elemento neutro}

  Moltiplicare qualsiasi numero per $1$ lascia il numero invariato:
  \[
    a \cdot 1 = 1 \cdot a = a
  \]
  L'\textbf{uno} è l'\textbf{elemento neutro} della moltiplicazione.
\end{regola}

\begin{esempio}
  \[
    37 \cdot 1 = 37 \qquad \qquad 1 \cdot 12 = 12
  \]
\end{esempio}

\begin{regola}
  \textbf{Proprietà dell'elemento assorbente}

  Moltiplicare qualsiasi numero per $0$ dà sempre $0$:
  \[
    a \cdot 0 = 0 \cdot a = 0
  \]
  Lo \textbf{zero} è l'\textbf{elemento assorbente} della moltiplicazione.
\end{regola}

\begin{esempio}
  \[
    999 \cdot 0 = 0 \qquad \qquad 0 \cdot 47 = 0
  \]
\end{esempio}

% ============================================================
%  SEZIONE 4 — DIVISIONE
% ============================================================
\section{La divisione}

\begin{definizione}
  La \textbf{divisione} è l'operazione che calcola quante volte un numero, detto \textbf{divisore}, è contenuto in un altro numero, detto \textbf{dividendo}. Il risultato si chiama \textbf{quoziente}.

  \[
    a : b = q
  \]

  dove $a$ è il dividendo, $b$ è il divisore e $q$ è il quoziente.

  \medskip
  Se la divisione non è esatta, si ottiene anche un \textbf{resto} $r$, con $0 \leq r < b$:
  \[
    a = b \cdot q + r
  \]
\end{definizione}

\subsection{Proprietà della divisione}

\begin{regola}
  \textbf{Proprietà invariantiva}

  Il quoziente non cambia se si moltiplicano (o si dividono) dividendo e divisore per uno stesso numero diverso da zero:
  \[
    a : b = (a \cdot c) : (b \cdot c) = (a : c) : (b : c) \qquad (c \neq 0)
  \]
\end{regola}

\begin{esempio}
  \[
    24 : 6 = 4
  \]
  \[
    (24 \cdot 3) : (6 \cdot 3) = 72 : 18 = 4
  \]
  \[
    (24 : 2) : (6 : 2) = 12 : 3 = 4
  \]
  Il quoziente è sempre $4$.
\end{esempio}

\begin{regola}
  \textbf{Elemento neutro parziale}

  Dividere qualsiasi numero per $1$ lascia il numero invariato:
  \[
    a : 1 = a
  \]
\end{regola}

\begin{esempio}
  \[
    45 : 1 = 45
  \]
\end{esempio}

\begin{attenzione}
  \textbf{Attenzione — Non si può dividere per zero!}

  L'operazione $a : 0$ \textbf{non è definita} per nessun valore di $a$.

  \medskip
  Dividere per zero non produce un risultato: è un'operazione che non ha senso matematico. In matematica scriviamo che $a : 0$ è \textbf{impossibile}.
\end{attenzione}

\begin{attenzione}
  \textbf{Attenzione — La divisione non è commutativa!}

  $a : b \neq b : a$ in generale.

  \medskip
  Ad esempio: $12 : 4 = 3$, ma $4 : 12 \neq 3$.
\end{attenzione}

% ============================================================
%  SEZIONE 5 — PRECEDENZA DELLE OPERAZIONI
% ============================================================
\section{La precedenza delle operazioni}

Quando un'espressione contiene più operazioni, occorre seguire un ordine preciso, detto \textbf{ordine di precedenza delle operazioni}.

\begin{regola}
  \textbf{Regola di precedenza}

  \begin{enumerate}
    \item Si calcolano prima le operazioni tra \textbf{parentesi}, partendo da quelle più interne.
    \item Poi si eseguono \textbf{moltiplicazioni e divisioni}, da sinistra verso destra.
    \item Infine si eseguono \textbf{addizioni e sottrazioni}, da sinistra verso destra.
  \end{enumerate}

  In sintesi: \textbf{parentesi} $\rightarrow$ \textbf{$\cdot$ e $:$} $\rightarrow$ \textbf{$+$ e $-$}
\end{regola}

\begin{esempio}
  Calcola: $\quad 3 + 4 \cdot 2 - (8 : 4)$

  \medskip
  \textbf{Passo 1 — Parentesi:}
  \[
    3 + 4 \cdot 2 - \mathbf{(8 : 4)} = 3 + 4 \cdot 2 - 2
  \]

  \textbf{Passo 2 — Moltiplicazione:}
  \[
    3 + \mathbf{4 \cdot 2} - 2 = 3 + 8 - 2
  \]

  \textbf{Passo 3 — Addizione e sottrazione (da sinistra a destra):}
  \[
    3 + 8 - 2 = 11 - 2 = 9
  \]

  \textbf{Risultato:} $9$
\end{esempio}

\begin{attenzione}
  \textbf{Errore comune — Calcolare da sinistra a destra senza rispettare la precedenza}

  \medskip
  \textbf{Sbagliato:} \quad $3 + 4 \cdot 2 = 7 \cdot 2 = 14$

  \textbf{Corretto:} \quad $3 + 4 \cdot 2 = 3 + 8 = 11$

  La moltiplicazione va eseguita \emph{prima} dell'addizione, anche se l'addizione appare prima nell'espressione scritta.
\end{attenzione}

% ============================================================
%  LO SAPEVI?
% ============================================================
\begin{sapevi}
  \textbf{Lo sapevi?}

  Il simbolo «$+$» per l'addizione e «$-$» per la sottrazione furono introdotti dal matematico tedesco Johannes Widmann nel 1489, nel suo libro \emph{Behende und hübsche Rechenung auff allen Kauffmanschafften}. Prima di allora, le operazioni venivano descritte a parole o con abbreviazioni latine come \emph{p} (per \emph{plus}) e \emph{m} (per \emph{minus}).
\end{sapevi}

% ============================================================
%  RIEPILOGO
% ============================================================
\section*{Riepilogo del capitolo}

\begin{center}
  \begin{tabular}{@{}llp{6cm}@{}}
    \toprule
    \textbf{Operazione} & \textbf{Proprietà}  & \textbf{Formula}                            \\
    \midrule
    Addizione           & Commutativa         & $a + b = b + a$                             \\
                        & Associativa         & $(a+b)+c = a+(b+c)$                         \\
                        & Elemento neutro     & $a + 0 = a$                                 \\
    \midrule
    Sottrazione         & Invariantiva        & $a - b = (a \pm c) - (b \pm c)$             \\
    \midrule
    Moltiplicazione     & Commutativa         & $a \cdot b = b \cdot a$                     \\
                        & Associativa         & $(a \cdot b) \cdot c = a \cdot (b \cdot c)$ \\
                        & Distributiva        & $a \cdot (b + c) = a \cdot b + a \cdot c$   \\
                        & Elemento neutro     & $a \cdot 1 = a$                             \\
                        & Elem.\ assorbente   & $a \cdot 0 = 0$                             \\
    \midrule
    Divisione           & Invariantiva        & $a : b = (a \cdot c) : (b \cdot c)$         \\
                        & Elem.\ neutro parz. & $a : 1 = a$                                 \\
                        & Divisione per 0     & $a : 0$ è \textbf{impossibile}              \\
    \midrule
    Precedenza          & Ordine              & Parentesi $\to$ $\cdot\,:\;$ $\to$ $+\,-$   \\
    \bottomrule
  \end{tabular}
\end{center}